\documentclass[../main.tex]{subfiles}

\begin{document}
\chapter{Introducción}
\label{C:introduccion}

Las máquinas Game \& Watch nacieron en 1980 y adquirieron una gran popularidad en la época. Muy resumidamente, la pantalla se dividía en segmentos sin intersección entre sí, que podían estar encendidos o apagados. A partir de esta idea básica se crearon multitud de juegos de muy diversas temáticas.

En este PFG comenzaríamos por analizar su historia, cómo surgieron, cómo funcionaban e incluso juegos posteriores para ordenador que adoptaron su misma filosofía, para proponer el diseño de nuestro propio motor para la creación de este tipo de juegos. Aunque lo importante será el diseño, para que cualquiera pueda, a partir de él, implementarlo en un lenguaje de su elección, se deberá implementar una versión sencilla del mismo.

El PFG no está orientado a una solución específica, sino que la persona que lo desarrolle podrá proponer cómo abordarlo tras analizar las distintas alternativas existentes (lo que permitirá situar el trabajo en el estado del arte).

El motor deberá proporcionar al desarrollador las acciones necesarias (gestión de pantalla, gestión de entrada, desarrollo de la partida, etc.) para implementar juegos de este tipo, tanto en pantalla simple como en pantalla doble.

El PFG podría incluir un ejemplo de uso del motor, es decir, un juego creado utilizando el motor a modo de demostración. También podría plantearse la posibilidad de crear un editor que permitiera generar juegos de forma sencilla.

\section{Motivación personal}

Como interesado en la historia de la informática, principalmente en las primeras máquinas electro mecánicas y las de lámparas de vacío, los Game \& Watch fueron las primeras máquinas populares de videojuegos, por lo que son parte de esa etapa inicial, donde se resolváin problemas con máquinas con muy pocos recursos.

Comencé a aprender a programar en los años 80, y hasta hoy ha sido mi profesión, aunque actualmente prefiero la arquitectura de soluciones al desarrollo, nunca le he perdido el gusto a programar como diversión, y este es un proyecto que será muy divertido.

Y como colofón,recuerdo muy bien haber pasado mucho tiempo jugando a una G\&W que compró mi padre en su día, por lo que hacer un emulador que me permita volver a jugar un poco a esta máquina es una idea muy atractiva para mi.

El decidirme por este TFG me permite por tanto divertirme programando, haciendo un emulador de una máquina que recuerdo con nostalgia, e investigar un poco más en esos aparatos y como funcionaban.

\section{Bibliografía}

No he localizado libros ni revistas que traten expresamente este tema, pero en Internet hay mucha información sobre estos aparatos, incluyendo simuladores on-line de los mismos, especialmente el proyecto MAME los incluye como uno de los sistemas que emula.

\end{document}